
\documentclass{memoir}

\usepackage{sourcecodepro}
\usepackage{sourcesanspro}
\usepackage{sourceserifpro}

\usepackage{fancyvrb}
\usepackage{fvextra}
\usepackage{xparse}

\usepackage[most]{tcolorbox}

% Detect whether we're in a footnote. This is used later to avoid \href rendering bugs.
\newif\ifinfootnote
\infootnotefalse

\let\oldfootnote\footnote
\renewcommand{\footnote}[1]{%
  \oldfootnote{\infootnotetrue#1\infootnotefalse}%
}

\let\oldfootnotetext\footnotetext
\renewcommand{\footnotetext}[1]{%
  \oldfootnotetext{\infootnotetrue#1\infootnotefalse}%
}

\usepackage[breaklinks=true, hyperfootnotes=false]{hyperref}

% Redefines \href to show footnotes with URLs instead. This works around a page breaking bug in
% hyperref and also makes the link useful on paper. If already in a footnote, the URL is
% in parentheses instead.
\let\oldhref\href
\RenewDocumentCommand{\href}{mm}{%
  \ifinfootnote%
    #2~(\url{#1})%
  \else%
    #2\footnote{\url{#1}}%
  \fi%
}

\usepackage[normalem]{ulem}
\newcommand{\coloredwave}[2]{\textcolor{#1}{\uwave{\textcolor{black}{#2}}}}
\usepackage{newunicodechar}

% Work around the fact that
% U+271D LATIN CROSS doesn't exist in
% DejaVu Sans Mono Oblique. \textup
% is fontspec for "upright, not italic/oblique".
\newunicodechar{✝}{\textup{✝}}
% Work around missing U+2011 (non-breaking hyphen) in Source Serif Pro
\newunicodechar{‑}{-}

\definecolor{errorColor}{HTML}{B91C1C}
\definecolor{infoColor}{HTML}{1E6BB8}
\definecolor{warningColor}{HTML}{D97706}
\newcommand{\errorDecorate}[1]{\coloredwave{errorColor}{#1}}
\newcommand{\infoDecorate}[1]{\coloredwave{infoColor}{#1}}
\newcommand{\warningDecorate}[1]{\coloredwave{warningColor}{#1}}
\DefineVerbatimEnvironment{LeanVerbatim}{Verbatim}
  {commandchars=\\\{\},fontsize=\small,breaklines=true}
\DefineVerbatimEnvironment{FileVerbatim}{Verbatim}{commandchars=\\\{\},fontsize=\small,breaklines=true,frame=single,framesep=2mm,numbers=left}
\CustomVerbatimCommand{\LeanVerb}{Verb}
  {commandchars=\\\{\},fontsize=\small,breaklines=true}
\CustomVerbatimCommand{\FileListingVerb}{Verb}
  {commandchars=\\\{\},fontsize=\small,frame=single,framesep=2mm, numbers=left}

%%% Trace messages
\newlength{\traceindent}
\setlength{\traceindent}{1.5em}

\newcommand{\expandedIndicator}{$\blacktriangledown$\hspace{2pt}}
\newcommand{\collapsedIndicator}{$\blacktriangleright$\hspace{2pt}}

\newenvironment{expandedtrace}[1]{%
  \par\noindent\expandedIndicator #1\par
  \advance\leftskip by \traceindent
}{%
}

\newenvironment{collapsedtrace}[1]{%
  \par\noindent\collapsedIndicator #1\par
  \advance\leftskip by \traceindent
}{%
}
%%% End Trace messages

\definecolor{bordercolor}{HTML}{98B2C0}
\definecolor{medgray}{HTML}{555555}
\newtcolorbox{docstringBox}[2][]{colback=white,
breakable,
colframe=bordercolor,
colbacktitle=white,
enhanced,
coltitle=medgray,
attach boxed title to top left={xshift=2mm,yshift=-2mm},
boxrule=0.4pt,
fonttitle=\sffamily\fontsize{6pt}{7pt}\selectfont,
boxed title style={top=-0.3mm,bottom=-0.3mm,left=-0.3mm,right=-0.3mm,boxrule=0.4pt},
title={#2},#1}


\makechapterstyle{lean}{%
\renewcommand*{\chaptitlefont}{\sffamily\HUGE}
\renewcommand*{\chapnumfont}{\chaptitlefont}
% allow for 99 chapters!
\settowidth{\chapindent}{\chapnumfont 999}
\renewcommand*{\printchaptername}{}
\renewcommand*{\chapternamenum}{}
\renewcommand*{\chapnumfont}{\chaptitlefont}
\renewcommand*{\printchapternum}{%
\noindent\llap{\makebox[\chapindent][l]{%
\chapnumfont \thechapter}}}
\renewcommand*{\afterchapternum}{}
}

\chapterstyle{lean}

\setsecheadstyle{\sffamily\bfseries\Large}
\setsubsecheadstyle{\sffamily\bfseries\large}
\setsubsubsecheadstyle{\sffamily\bfseries}

\renewcommand{\cftchapterfont}{\normalfont\sffamily}
\renewcommand{\cftsectionfont}{\normalfont\sffamily}
\renewcommand{\cftchapterpagefont}{\normalfont\sffamily}
\renewcommand{\cftsectionpagefont}{\normalfont\sffamily}
\setmonofont{DejaVu Sans Mono}

\definecolor{textgray}{HTML}{2C2C2C}
\definecolor{codecomment}{rgb}{0.497495, 0.497587, 0.497464}
\usepackage[a4paper,margin=1in]{geometry}
\usepackage{parskip}
\usepackage{amsfonts}
\AtBeginDocument{\color{textgray}}
\renewcommand*{\chaptitlefont}{\sffamily\HUGE\color{textgray}}
\renewcommand*{\chapnumfont}{\chaptitlefont}
\setsecheadstyle{\sffamily\bfseries\Large\color{textgray}}
\setsubsecheadstyle{\sffamily\bfseries\large\color{textgray}}
\setsubsubsecheadstyle{\sffamily\bfseries\color{textgray}}
\setmonofont{CMU Typewriter Text}
\newfontfamily{\PLFASymbolFont}{DejaVu Sans Mono}[Scale=0.6]
\newcommand{\PLFANat}{\ensuremath{\mathbb{N}}}
\newcommand{\PLFAForall}{\ensuremath{\forall}}
\newcommand{\PLFATo}{\ensuremath{\to}}
\newcommand{\PLFAEquiv}{\ensuremath{\equiv}}
\newcommand{\PLFATurnstile}{\ensuremath{\vdash}}
\newcommand{\PLFASymbol}[1]{{\PLFASymbolFont #1}}

% Keep the PDF viewer layout book-like, but stop memoir from alternating
% recto/verso headers and blank-page behavior.
\makeatletter
\@twosidefalse
\@mparswitchfalse
\makeatother
\let\cleardoublepage\clearpage
\makeevenhead{headings}{\slshape\color{textgray}\rightmark}{}{\color{textgray}\thepage}
\makeoddhead{headings}{\slshape\color{textgray}\rightmark}{}{\color{textgray}\thepage}

% Verso emits fenced blocks as verbatim/LeanVerbatim. Tighten the display
% spacing so code blocks do not leave a large hole before the next paragraph.
\newcommand{\PLFACodeSize}{\normalsize}
\newcommand{\PLFAFirstOfOne}[1]{#1}
\newcommand{\PLFAComment}[1]{\textcolor{codecomment}{#1}}
\RecustomVerbatimEnvironment{verbatim}{Verbatim}
  {commandchars=\\\{\},formatcom=\ttfamily,fontsize=\PLFACodeSize,breaklines=true,baselinestretch=0.92,
   listparameters={\setlength{\topsep}{0.25\baselineskip}\setlength{\partopsep}{0pt}\setlength{\parsep}{0pt}}}
\RecustomVerbatimEnvironment{LeanVerbatim}{Verbatim}
  {commandchars=\\\{\},formatcom=\ttfamily\let\textit\PLFAFirstOfOne,fontsize=\PLFACodeSize,breaklines=true,baselinestretch=0.92,
   listparameters={\setlength{\topsep}{0.25\baselineskip}\setlength{\partopsep}{0pt}\setlength{\parsep}{0pt}}}

\title{\sffamily PLFA in Lean}
\author{\sffamily Daniel}
\date{\sffamily }

\begin{document}

\frontmatter

\begin{titlingpage}
\maketitle
\end{titlingpage}

\tableofcontents

\par

\cleardoublepage
\clearpage

\par

\mainmatter
\chapter{Part 1}\label{PLFA--in--Lean----Part--1}

\par
This part collects the early material on natural numbers and induction.

\par

\section{Naturals}\label{PLFA--in--Lean----Part--1----Naturals}

\par
Most are familiar with the natural numbers:

\par
\begin{verbatim}
0
1
2
3
...
\end{verbatim}

\par
\par\noindent and so on. The set of natural numbers \LeanVerb|\PLFANat{}| is inifite and we say that \LeanVerb|0|, \LeanVerb|1|,
\LeanVerb|2|, \LeanVerb|3|, and are \textbf{values} of type \LeanVerb|\PLFANat{}| (otherwise indicated with: \LeanVerb|0 : \PLFANat{}|,
\LeanVerb|1 : \PLFANat{}|,  etc.).

\par
Although there are infinitely many naturals, yet we can write down its
definition in just a few lines. Here it is as a pair of inference rules:

\par
\begin{verbatim}
--------
zero : \PLFANat{}

m : \PLFANat{}
--------
suc m : \PLFANat{}
\end{verbatim}

\par
And here is the definition in Lean:

\par
\begin{LeanVerbatim}
\textbf{inductive} \PLFANat{} : Type \textbf{where}
    \symbol{124} zero : \PLFANat{}
    \symbol{124} suc  : \PLFANat{} -> \PLFANat{}
\end{LeanVerbatim}

\par
Here \LeanVerb|\PLFANat{}| is the name of the datatype we are defining, and \LeanVerb|zero| and \LeanVerb|suc|
(short for successor) are the \textbf{constructors} of the datatype.

\par
Both definitions above tell us two things:

\begin{enumerate}
\item base case: \LeanVerb|zero| is a natural number
\item inductive case: if \LeanVerb|m| is a natural number, then \LeanVerb|suc m| is also a natural
number.

\end{enumerate}

\par
Further, these two rules give the \textbf{only} ways of creating natural numbers.
Hence, the possible natural numbers are:

\par
\begin{verbatim}
.zero
.suc .zero
.suc (.suc .zero)
...
\end{verbatim}

\par
We write \LeanVerb|0| as the shorthand for the \LeanVerb|zero| constructor and \LeanVerb|1| as shorthand
for \LeanVerb|suc zero|, the successor of zero, that is the number that comes after \LeanVerb|0|.

\par

\subsection{Exercise: \texttt{seven} (practice)}\label{PLFA--in--Lean----Part--1----Naturals----Exercise-005F-005F-005F----seven-----005FLPAR-005Fpractice-005FRPAR-005F}

\par
Write out \LeanVerb|7| in longhand.

\par
\textbf{potential sol:}

\par
\begin{LeanVerbatim}
\textbf{def} seven : \PLFANat{} :=
  .suc $ .suc $ .suc $ .suc $ .suc $ .suc $ .suc .zero
\end{LeanVerbatim}

\par

\subsection{Unpacking the inference rules}\label{PLFA--in--Lean----Part--1----Naturals----Unpacking--the--inference--rules}

\par
Let's unpack the Lean definition. the keyword \LeanVerb|inductive| tells us this is an
inductive definition, that is, we are defining a new datatype with constructors.
The phrase:

\par
\begin{verbatim}
\PLFANat{} : Type
\end{verbatim}

\par
tells us that \LeanVerb|\PLFANat{}| is the name of the new data, type and that it is a \LeanVerb|Type|,
which is the way in Lean of saying that it is a type.

\par
The keyword \LeanVerb|where| is the name of the new datatype, and that is a \LeanVerb|Type|,
which is the way in Lean of saying that it is a type. The keyword \LeanVerb|where|
separates the declaration of the data from the declaration of its constructors.
Each constructor is a declared on a separate line, which is indented to
indicate that it belongs to the corresponding \LeanVerb|data| declaration. The lines:

\par
\begin{verbatim}
zero : \PLFANat{}
suc  : \PLFANat{} -> \PLFANat{}
\end{verbatim}

\par
give \emph{signatures} specifying the types of the constructors \LeanVerb|zero| and \LeanVerb|suc|.
They tell us \LeanVerb|zero| is a natural number and that \LeanVerb|suc| takes a natural number as
an argument and returns a natural number.

\par
You may have noticed that \LeanVerb|\PLFANat{}| and \LeanVerb|\PLFATo{}| don't appear on your keyboard. They are
symbols in \emph{unicode}. At the end of each chapter is a list of all unicode
symbols introduced in the chapter, including instruction on how to insert them
into the editor (assuming you're using zed or intellij).

\par

\subsection{Lean4 Natural Literal Syntax}\label{PLFA--in--Lean----Part--1----Naturals----Lean4--Natural--Literal--Syntax}

\par
Lean4 already defines its own inductive type for the natural numbers and also
supports integer literal syntax. We can hook into that sytax and relate it to
our own \LeanVerb|\PLFANat{}| type with the following conversion function and typeclass instance
implementation:

\par
\begin{LeanVerbatim}
\textbf{def} convert : _root_.Nat -> \PLFANat{}
    \symbol{124} _root_.Nat.zero   => .zero
    \symbol{124} _root_.Nat.succ \textit{n} => .suc (convert \textit{n})

\textbf{instance} (\textit{n} : _root_.Nat) : OfNat \PLFANat{} \textit{n} \textbf{where}
    ofNat := convert \textit{n}
\end{LeanVerbatim}

\par

\subsection{Operations on naturals are recursive functions}\label{PLFA--in--Lean----Part--1----Naturals----Operations--on--naturals--are--recursive--functions}

\par
Now that we have the natural numbers, what can we do with them? For instance,
can we define arithmetic operations such as addition and multiplication?

\par
It turns out, each of these can be described as \emph{recursive} functions. Here is
the definition of addition for the \LeanVerb|\PLFANat{}| type in lean4:

\par
\begin{LeanVerbatim}
\textbf{def} plus : \PLFANat{} -> \PLFANat{} -> \PLFANat{}
    \symbol{124} .zero , \textit{n}    => \textit{n}
    \symbol{124} (.suc \textit{m}) , \textit{n} => .suc (plus \textit{m} \textit{n})

\textbf{syntax}:65 (\textbf{priority} := \textbf{high}) term:66 " + " term:65 : term

\textbf{macro_rules}
  \symbol{124} `($\textit{a} + $\textit{b}) => `(plus $\textit{a} $\textit{b})
\end{LeanVerbatim}

\par
The first line of \LeanVerb|plus| gives its type: \LeanVerb|\PLFANat{} -> \PLFANat{} -> \PLFANat{}|, which indicates that
the \LeanVerb|plus| function accepts two naturals and returns a natural. The infix
notation allows us to write plus using the usual infix \LeanVerb|+| notation. The
\LeanVerb|priority := high| bit makes it more likely that uses of
\LeanVerb|+| will get resolved to our custom \LeanVerb|plus| function (for our own natural
number type (\LeanVerb|\PLFANat{}|) based on the surrounding context).

\par
If we write zero as \LeanVerb|0| and \LeanVerb|suc m| as \LeanVerb|1 + m|, the definition turns into two familiar equations:

\par
\begin{verbatim}
 0       + n  =  n
 (1 + m) + n  =  1 + (m + n)
\end{verbatim}

\par
\par\noindent The first follows because zero is an identity for addition, and the second because addition is associative.
In its most general form, associativity is written\begin{verbatim}
(m + n) + p = m + (n + p)
\end{verbatim}
meaning the location of the parentheses is irrelevant.

\par
An example of how we can use these constructors (better more illustrative
example second I think):

\par
\begin{LeanVerbatim}
\textbf{example} : 3 + 2 = 5 :=
  \textbf{calc}
      3 + 2
      = .suc (2 + 2)                 := rfl
      _ = .suc (.suc (1 + 2))        := rfl
      _ = .suc (.suc (.suc (0 + 2))) := rfl
      _ = .suc (.suc (.suc 2))       := rfl
      _ = 5                          := rfl
\end{LeanVerbatim}

\par
Note we could've expressed the theorem like so if we wanted to
make it more clear in the local context that these literals are formed via our
\LeanVerb|\PLFANat{}| type like so: \LeanVerb|example : 3 + 2 = (5 : \PLFANat{})|

\par
Here's another version closer to the original PLFA text:

\par
\begin{LeanVerbatim}
\textbf{section}
\textbf{open} \PLFANat{}  \PLFAComment{-- so we don't need dots in front of each ctor}
\textbf{example} : 2 + 3 = 5 :=
  \textbf{calc}
      2 + 3
      = (suc (suc zero)) + (suc (suc (suc zero)))   := rfl
      _ = suc ((suc zero) + (suc (suc (suc zero)))) := rfl
      _ = suc (suc (zero + (suc (suc (suc zero))))) := rfl
      _ = suc (suc (suc (suc (suc zero))))          := rfl
      _ = 5                                         := rfl
\textbf{end}
\end{LeanVerbatim}

\par

\subsection{Multiplication}\label{PLFA--in--Lean----Part--1----Naturals----Multiplication}

\par
Once we have defined addition, we can define multiplication as repeated
addition:

\par
\begin{LeanVerbatim}
\textbf{def} mult : \PLFANat{} -> \PLFANat{} -> \PLFANat{}
    \symbol{124} .zero, \warningDecorate{\textit{n}}      => .zero
    \symbol{124} (.suc \textit{m}), \textit{n}   => \textit{n} + (mult \textit{m} \textit{n})

\PLFAComment{-- infix syntax rules for infix syntax:}
\textbf{syntax}:70 (\textbf{priority} := \textbf{high}) term:70 " * " term:71 : term

\textbf{macro_rules}
  \symbol{124} `($\textit{a} * $\textit{b}) => `(mult $\textit{a} $\textit{b})
\end{LeanVerbatim}

\par
\par\noindent Computing \LeanVerb|m * n| returns the sum of \LeanVerb|m| copies of \LeanVerb|n|. Again, rewriting turns
the definition into two familiar equations:

\par
\begin{verbatim}
0       * n  =  0
(1 + m) * n  =  n + (m * n)
\end{verbatim}

\par

\subsection{Exercise \texttt{*-example} (practice)}\label{PLFA--in--Lean----Part--1----Naturals----Exercise-----005F-005F-005F--example-----005FLPAR-005Fpractice-005FRPAR-005F}

\par
Compute \LeanVerb|3 * 4| writing out your reasoning as a chain of equations using the
equations for \LeanVerb|*|. You do not need to step through the evaluation of \LeanVerb|+|

\par
\textbf{potential sol:}

\par
\begin{LeanVerbatim}
\textbf{example} : 3 * 4 = 12 :=
    \textbf{calc}
        3 * 4
      = 4 + (2 * 4)              := rfl
    _ = 4 + (4 + (1 * 4))        := rfl
    _ = 4 + (4 + (4 + (0 * 4)))  := rfl
    _ = 4 + (4 + (4 + 0))        := rfl
    _ = 4 + 4 + 4                := rfl
    _ = 12                       := rfl
\end{LeanVerbatim}

\par

\subsection{Exercise exponentiation (recommended)}\label{PLFA--in--Lean----Part--1----Naturals----Exercise--exponentiation---005FLPAR-005Frecommended-005FRPAR-005F}

\par
Define exponentiation, which is given by the following equations:

\par
\begin{verbatim}
m ^ 0           = 1
m ^ (1 + n)     = m * (m^n)
\end{verbatim}

\par
\textbf{potential sol:}

\par
\begin{LeanVerbatim}
\textbf{def} exponent : \PLFANat{} -> \PLFANat{} -> \PLFANat{}
    \symbol{124} \warningDecorate{\textit{m}} , 0         => 1
    \symbol{124} \textit{m} , (.suc \textit{n})  => \textit{m} * (exponent \textit{m} \textit{n})

\textbf{syntax}:80 (\textbf{priority} := \textbf{high}) term:81 " ^ " term:80 : term

\textbf{macro_rules}
  \symbol{124} `($\textit{a} ^ $\textit{b}) => `(exponent $\textit{a} $\textit{b})
\end{LeanVerbatim}

\par
Check that \LeanVerb|3 ^ 4| is \LeanVerb|81|:

\par
\begin{LeanVerbatim}
\infoDecorate{\textbf{\symbol{35}check}} (3 ^ 4 = 81)
\end{LeanVerbatim}

\par

\subsection{Monus}\label{PLFA--in--Lean----Part--1----Naturals----Monus}

\par
We can also define subtraction for natural numbers. Since there are no negative
natural numbers, if we subtract a larger number from a smaller number we will
take the result to be zero. This adaption of subtraction to naturals is called
monus (a twist on \emph{minus}).

\par
Monus is our first use of a definition that uses pattern matching against both
arguments:

\par
\begin{LeanVerbatim}
\textbf{def} monus : \PLFANat{} -> \PLFANat{} -> \PLFANat{}
    \symbol{124} \textit{m}, .zero              => \textit{m}
    \symbol{124} .zero, (.suc \warningDecorate{\textit{n}})       => .zero
    \symbol{124} (.suc \textit{m}), (.suc \textit{n})    => monus \textit{m} \textit{n}

\textbf{syntax}:80 (\textbf{priority} := \textbf{high}) term:81 " \PLFASymbol{∸} " term:80 : term

\textbf{macro_rules}
  \symbol{124} `($\textit{a} \PLFASymbol{∸} $\textit{b}) => `(monus $\textit{a} $\textit{b})
\end{LeanVerbatim}

\par
\par\noindent where lean performs case analysis to ensure that all the cases are covered.

\par

\begin{itemize}
\item consider the second argument:
\item \begin{itemize}
\item if it is \LeanVerb|zero|, then the first equation applies

\end{itemize}

\item \begin{itemize}
\item if it is \LeanVerb|suc n|, then consider the first argument

\end{itemize}

\item \begin{itemize}
\item a) if it is \LeanVerb|zero|, then the second equation applies

\end{itemize}

\item \begin{itemize}
\item b) if it is \LeanVerb|suc m|, then the third equation applies

\end{itemize}

\par

\end{itemize}

\par
Lean will raise an error if all the cases are not covered. As with addition
and multiplication, the recursive definition is well founded because monus on
bigger numbers is defined in terms of monus on smaller numbers.

\par
For example, let's subtract two from three:

\par
\begin{LeanVerbatim}
\textbf{example} : 3 \PLFASymbol{∸} 2 = 1 :=
    \textbf{calc}
        2 \PLFASymbol{∸} 1
      = 1 \PLFASymbol{∸} 0 := rfl
    _ = 1     := rfl
\end{LeanVerbatim}

\par
We did not use the second equation at all, but it will be required if we try to
subtract a larger number from a smaller one:

\par
\begin{LeanVerbatim}
\textbf{example} : 2 \PLFASymbol{∸} 3 = 0 :=
    \textbf{calc}
        1 \PLFASymbol{∸} 2
    _ = 0 \PLFASymbol{∸} 1        := rfl \PLFAComment{-- desugars to:}
    _ = 0 \PLFASymbol{∸} (.suc 0) := rfl
    _ = 0            := rfl
\end{LeanVerbatim}

\par

\subsection{Currying}\label{PLFA--in--Lean----Part--1----Naturals----Currying}

\par
We have chosen to represent a function of two arguments in terms of a function
of the first argument that returns a function of the second argument. This trick
goes by the name \emph{currying}.

\par
Lean4, like other functional languages such as Haskell and ML, is designed to
make currying easy to use. Function arrows associate to the right and application
associates to the left:

\par
\LeanVerb|\PLFANat{} -> \PLFANat{} -> \PLFANat{}| associates like so: \LeanVerb|\PLFANat{} -> (\PLFANat{} -> \PLFANat{})|

\par
\par\noindent and

\par
\LeanVerb|plus 2 3| stands for \LeanVerb|(plus 2) 3|

\par
\par\noindent The term \LeanVerb|plus 2| by itself stands for the function that adds two to its
argument, hence applying it to three yields five.

\par
Currying is named for Haskell B. Curry (a Penn State professor actually:
https://www.hmdb.org/m.asp?m=134735;
\href{https://mathshistory.st-andrews.ac.uk/Biographies/Curry/}{bio})

\par

\subsection{Story of creation, revisited}\label{PLFA--in--Lean----Part--1----Naturals----Story--of--creation-005F-005F-005F--revisited}

\par
Just as our inductive definition defines the naturals in terms of the naturals,
so does our recursive definition define addition in terms of addition.

\par
Again we resort to a creation story, where this time we are concerned with
judgments about addition:

\par
\begin{verbatim}
-- In the beginning, we know nothing about addition.
\end{verbatim}

\par
Now, we apply the rules to all the judgment we know about. The base case tells
us that \LeanVerb|.zero + n = n| for every natural \LeanVerb|n|, so we add all those equations.
The inductive case tells us that if \LeanVerb|m + n = p| (on the day before today) then
\LeanVerb|suc m + n = suc p| (today). We didn't know any equations about addition before
today, so that rule doesn't give us any new equations:

\par
\begin{verbatim}
-- On the first day, we know about addition of 0.
0 + 0 = 0     0 + 1 = 1    0 + 2 = 2     ...
\end{verbatim}

\par
Then we repeat the process, so on the next day we know about all the equations
from the day before, plus any equations added by the rules. The base case tells
us nothing new, but now the inductive case adds more equations:

\par
\begin{verbatim}
-- On the second day, we know about addition of 0 and 1.
0 + 0 = 0     0 + 1 = 1     0 + 2 = 2     0 + 3 = 3     ...
1 + 0 = 1     1 + 1 = 2     1 + 2 = 3     1 + 3 = 4     ...
\end{verbatim}

\par
And we repeat the process again:

\par
\begin{verbatim}
-- On the third day, we know about addition of 0, 1, and 2.
0 + 0 = 0     0 + 1 = 1     0 + 2 = 2     0 + 3 = 3     ...
1 + 0 = 1     1 + 1 = 2     1 + 2 = 3     1 + 3 = 4     ...
2 + 0 = 2     2 + 1 = 3     2 + 2 = 4     2 + 3 = 5     ...
\end{verbatim}

\par
You've got the hang of it by now:

\par
\begin{verbatim}
-- On the fourth day, we know about addition of 0, 1, 2, and 3.
0 + 0 = 0     0 + 1 = 1     0 + 2 = 2     0 + 3 = 3     ...
1 + 0 = 1     1 + 1 = 2     1 + 2 = 3     1 + 3 = 4     ...
2 + 0 = 2     2 + 1 = 3     2 + 2 = 4     2 + 3 = 5     ...
3 + 0 = 3     3 + 1 = 4     3 + 2 = 5     3 + 3 = 6     ...
\end{verbatim}

\par
The process continues. On the m'th day we will know all the equations where the
first number is less than \LeanVerb|m|.

\par
As we can see, the reasoning that justifies inductive and recursive definitions
is quite similar. They might be considered two sides of the same coin.

\par

\subsection{The story of creation, finitely}\label{PLFA--in--Lean----Part--1----Naturals----The--story--of--creation-005F-005F-005F--finitely}

\par
The above story was told in a stratified way. First, we create the infinite set
of naturals. We take that set as given when creating instances of addition, so
even on day one we have an infinite set of instances.

\par
Instead, we could choose to create both the naturals and the instances of
addition at the same time. Then on any day there would be only a finite set of
instances:

\par
\begin{verbatim}
-- In the beginning, we know nothing.
\end{verbatim}

\par
Now, we apply the rules to all the judgment we know about. Only the base case
for naturals applies:

\par
\begin{verbatim}
-- On the first day, we know zero.
0 : \PLFANat{}
\end{verbatim}

\par
Again, we apply all the rules we know. This gives us a new natural, and our
first equation about addition.

\par
\begin{verbatim}
-- On the second day, we know one and all sums that yield zero.
0 : \PLFANat{}
1 : \PLFANat{}    0 + 0 = 0
\end{verbatim}

\par
Then we repeat the process. We get one more equation about addition from the
base case, and also get an equation from the inductive case, applied to
equation of the previous day:

\par
\begin{verbatim}
-- On the third day, we know two and all sums that yield one.
0 : \PLFANat{}
1 : \PLFANat{}    0 + 0 = 0
2 : \PLFANat{}    0 + 1 = 1   1 + 0 = 1
\end{verbatim}

\par
\par\noindent You've got the hang of it by now:

\par
\begin{verbatim}
-- On the fourth day, we know three and all sums that yield two.
0 : \PLFANat{}
1 : \PLFANat{}    0 + 0 = 0
2 : \PLFANat{}    0 + 1 = 1   1 + 0 = 1
3 : \PLFANat{}    0 + 2 = 2   1 + 1 = 2    2 + 0 = 2
\end{verbatim}

\par
On the n'th day there will be n distinct natural numbers, and \LeanVerb|n \PLFASymbol{×} (n-1) / 2|
equations about addition. The number \LeanVerb|n| and all equations for addition of
numbers less than \LeanVerb|n| first appear by day \LeanVerb|n+1|. This gives an entirely
finitist view of infinite sets of data and equations relating the data.

\par

\subsection{Exercise \texttt{Bin} (stretch)}\label{PLFA--in--Lean----Part--1----Naturals----Exercise----Bin-----005FLPAR-005Fstretch-005FRPAR-005F}

\par
A more efficient representation of natural numbers uses a binary rather than a
unary system. We represent a number as a bitstring:

\par
\begin{LeanVerbatim}
\textbf{inductive} Bin : Type \textbf{where}
    \symbol{124} empty  : Bin
    \symbol{124} t_zero : Bin -> Bin
    \symbol{124} t_one  : Bin -> Bin
    \textbf{deriving} Repr

\textbf{syntax}:80 (\textbf{priority} := \textbf{high}) "\PLFASymbol{⟨}\PLFASymbol{⟩}" : term
\textbf{syntax}:81 term:80 "O" : term
\textbf{syntax}:81 term:80 "I" : term
\textbf{macro_rules}
  \symbol{124} `(\PLFASymbol{⟨}\PLFASymbol{⟩})        => `(Bin.empty)
  \symbol{124} `($\textit{a}:\textbf{term} \textbf{O}) => `(Bin.t_zero $\textit{a})
  \symbol{124} `($\textit{a}:\textbf{term} \textbf{I}) => `(Bin.t_one $\textit{a})
\end{LeanVerbatim}

\par
\par\noindent So the bitstring:

\par
\begin{verbatim}
1011
\end{verbatim}

\par
\par\noindent would stand for the number eleven, encoded in our meta-introduced notation like so:\begin{LeanVerbatim}
\infoDecorate{\textbf{\symbol{35}eval}} \PLFASymbol{⟨}\PLFASymbol{⟩} \textbf{I} \textbf{O} \textbf{I} \textbf{I}
\end{LeanVerbatim}
written desugared:\begin{verbatim}
Bin.t_one (Bin.t_one (Bin.t_zero (Bin.t_one (Bin.empty))))
\end{verbatim}
Representations are not unique due to leading zeroes. Hence eleven can also be
represented by \LeanVerb|0001011|, encoded as:\begin{LeanVerbatim}
\infoDecorate{\textbf{\symbol{35}eval}} \PLFASymbol{⟨}\PLFASymbol{⟩} \textbf{O} \textbf{O} \textbf{I} \textbf{O} \textbf{I} \textbf{I}
\end{LeanVerbatim}
Define a function:\begin{verbatim}
def inc : Bin -> Bin
\end{verbatim}
that converts a bitstring to the bitstring for the next higher number.

\par

\clearpage

\par

\section{Induction: Proof by Induction}\label{PLFA--in--Lean----Part--1----Induction-005F-005F-005F--Proof--by--Induction}

\par
Now that we've defined the naturals and operations on them, our next step is to
learn how to prove properties that they satisfy. As hinted by their name,
properties of \emph{inductive datatypes} are proved by \emph{induction}.

\par

\subsection{Properties of operators}\label{PLFA--in--Lean----Part--1----Induction-005F-005F-005F--Proof--by--Induction----Properties--of--operators}

\par
Operations pop up all the time, and mathematicians have agreed on names for some
of the most common properties.

\par
\par\noindent \begin{itemize}
\item \emph{Identity}. Operator \LeanVerb|+| has left identity \LeanVerb|0| if \LeanVerb|0 + n = n|, and right
identity if \LeanVerb|n + 0 = n|, for all \LeanVerb|n|. A value that is both a left and right
identity is just called an identity. Identity is also sometimes called \emph{unit}.
\item \emph{Associativity}. Operator \LeanVerb|+| is associative if the location of the parentheses
does not matter: \LeanVerb|(m + n) + p = m + (n + p)|
\item \emph{Commutativity}. Operator \LeanVerb|+| is commutative if the order of arguments does not
matter: \LeanVerb|m + n = n + m|, for all \LeanVerb|m| and \LeanVerb|n|.
\item \emph{Distributivity}. Operator \LeanVerb|*| distributes over operator \LeanVerb|+| from the left if
\LeanVerb|m * (p + q) = (m * p) + (m * q)| for all \LeanVerb|m|, \LeanVerb|p|, and \LeanVerb|q|, and from the right
if \LeanVerb|(m + n) * p = (m * p) + (n * p)| for all \LeanVerb|m|, \LeanVerb|n|, and \LeanVerb|p|.

\end{itemize}
Addition has identity \LeanVerb|0| and multiplication has identity \LeanVerb|1|; addition and
multiplication are both associative and commutative; and multiplication
distributes over addition.

\par

\subsection{Associativity}\label{PLFA--in--Lean----Part--1----Induction-005F-005F-005F--Proof--by--Induction----Associativity}

\par
One property of addition is that it is associative, that is, that the location
of the parentheses does not matter:

\par
\par\noindent \begin{verbatim}
(m + n) + p \PLFAEquiv{} m + (n + p)
\end{verbatim}
Here \LeanVerb|m|, \LeanVerb|n|, and \LeanVerb|p| are variables that range over all natural numbers.We can test the proposition by choosing specific numbers for the three variables:\begin{verbatim}
example : (3 + 4) + 5 = 3 + (4 + 5) := by
  calc
    (3 + 4) + 5 = 7 + 5 := by rfl
    _ = 12              := by rfl
    _ = 3 + 9           := by rfl
    _ = 3 + (4 + 5)     := by rfl
\end{verbatim}
Here we have displayed the computation as a chain of equations, one term to a
line. It is often easiest to read such chains from the top down until one
reaches the simplest term (in this case, \LeanVerb|12|), and then from the bottom up
until one reaches the same term.The test reveals that associativity is perhaps not as obvious as it first
appears. Why should \LeanVerb|7 + 5| be the same as \LeanVerb|3 + 9|? We might want to gather
more evidence by testing the proposition with other numbers. But since there
are infinitely many naturals, testing can never be complete. Is there any way
we can be sure that associativity holds for all the natural numbers?

\par
\par\noindent The answer is yes! We can prove a property holds for all naturals using proof by induction.

\par

\subsection{Proof by induction}\label{PLFA--in--Lean----Part--1----Induction-005F-005F-005F--Proof--by--Induction----Proof--by--induction}

\par
Recall that the definition of natural numbers consists of a base case which
tells us that \LeanVerb|zero| is a natural, and an inductive case which tells us that if
\LeanVerb|m| is a natural then \LeanVerb|suc m| is also a natural.

\par
\par\noindent Proof by induction follows the structure of this definition. To prove a property
of natural numbers by induction, we need to prove two cases. First is the base
case, where we show the property holds for \LeanVerb|zero|. Second is the inductive case,
where we assume the property holds for an arbitrary natural \LeanVerb|m| (we call this
the inductive hypothesis), and then show that the property must also hold for
\LeanVerb|suc m|.

\par
If we write \LeanVerb|P m| for a property of \LeanVerb|m|, then what we need to demonstrate are
the following two inference rules:

\par
\par\noindent \begin{verbatim}
------
P .zero

P m
---------
P (.suc m)
\end{verbatim}

\par
Let's unpack these rules. The first rule is the base case, and requires us to
show that property \LeanVerb|P| holds for \LeanVerb|zero|. The second rule is the inductive case,
and requires us to show that if we assume the inductive hypothesis, namely that
\LeanVerb|P| holds for \LeanVerb|m|, then it follows that \LeanVerb|P| also holds for \LeanVerb|suc m|.

\par
Why does this work? Again, it can be explained by a creation story. To start
with, we know no properties:

\par
\par\noindent \begin{verbatim}
-- in the beginning, no properties are known
\end{verbatim}
Now, we apply the two rules to all the properties we know about. The base case
tells us that \LeanVerb|P zero| holds, so we add it to the set of known properties. The
inductive case tells us that if \LeanVerb|P m| holds on the day before today, then
\LeanVerb|P (suc m)| also holds today. We didn’t know about any properties before today,
so the inductive case doesn’t apply:\begin{verbatim}
-- on the first day, one property is known
P .zero
\end{verbatim}
Then we repeat the process, so on the next day we know about all the properties
from the day before, plus any properties added by the rules. The base case
tells us that \LeanVerb|P zero| holds, but we already knew that. Now the inductive case
tells us that since \LeanVerb|P zero| held yesterday, then \LeanVerb|P (suc zero)| holds today:\begin{verbatim}
-- on the second day, two properties are known
P .zero
P (.suc .zero)
\end{verbatim}
And we repeat the process again. Now the inductive case tells us that since
\LeanVerb|P zero| and \LeanVerb|P (suc zero)| both hold, then \LeanVerb|P (suc zero)| and
\LeanVerb|P (suc (suc zero))| also hold. We already knew about the first of these, but
the second is new:\begin{verbatim}
-- on the third day, three properties are known
P .zero
P (.suc .zero)
P (.suc (.suc .zero))
\end{verbatim}
You've got the hang of it by now:\begin{verbatim}
-- on the fourth day, four properties are known
P  .zero
P (.suc .zero)
P (.suc (.suc .zero))
P (.suc (.suc (.suc .zero)))
\end{verbatim}
The process continues. On the \LeanVerb|n|-th day there will be \LeanVerb|n| distinct properties
that hold. The property of every natural number will appear on some given day.
In particular, the property \LeanVerb|P n| first appears on day \LeanVerb|n + 1|.

\par

\subsection{Our first proof: associativity}\label{PLFA--in--Lean----Part--1----Induction-005F-005F-005F--Proof--by--Induction----Our--first--proof-005F-005F-005F--associativity}

\par
\par\noindent To prove associativity, we take \LeanVerb|P m| to be the property:\begin{verbatim}
plus (plus m n) p = plus m (plus n p)`
\end{verbatim}
Here \LeanVerb|n| and \LeanVerb|p| are arbitrary natural numbers, so if we can show the equation
holds for all \LeanVerb|m| it will also hold for all \LeanVerb|n| and \LeanVerb|p|. The appropriate
instances of the inference rules are:\begin{verbatim}
-------------------------------
plus (plus .zero n) p = plus .zero (plus n p)

plus (plus m n) p = plus m (plus n p)
---------------------------------
(.suc (plus (plus m n) p) = .suc (plus m (plus n p))
\end{verbatim}
If we can demonstrate both of these, then associativity of addition follows by
induction.Here is the proposition's statement and proof:\begin{LeanVerbatim}
\textbf{theorem} plus_assoc :
  \PLFAForall{} (\textit{m} \textit{n} \textit{p} : \PLFANat{}), (\textit{m} + \textit{n}) + \textit{p} = \textit{m} + (\textit{n} + \textit{p}) := \textbf{by}
    \textbf{intro} \textit{m} \textit{n} \textit{p}
    \textbf{induction} \textit{m} \textbf{with}
    \PLFAComment{--goal: \PLFAForall{} m n p : \PLFANat{},}
    \PLFAComment{--      \PLFATurnstile{} plus (plus m n) p = plus m (plus n p)}
    \symbol{124} zero =>
        \textbf{calc}
            (.zero + \textit{n}) + \textit{p} = (\textit{n} + \textit{p})   := \textbf{by} \textbf{rfl}
          _ = .zero + (\textit{n} + \textit{p})           := \textbf{by} \textbf{rfl}
    \symbol{124} suc \textit{m} \textit{ih} =>
        \PLFAComment{-- ih : plus (plus m n) p = plus m (plus n p)}
        \PLFAComment{-- ind. case goal:}
        \PLFAComment{--  \PLFATurnstile{} plus (plus m.suc n) p = plus m.suc (plus n p)}
        \textbf{calc}
            plus (plus (.suc \textit{m}) \textit{n}) \textit{p}
            = plus (.suc (plus \textit{m} \textit{n})) \textit{p} := \textbf{by} \textbf{rfl}
          _ = .suc (plus \textit{m} (plus \textit{n} \textit{p})) := \textbf{by}
              \textbf{exact} congrArg \PLFANat{}.suc (\textit{ih})
          _ = .suc (plus \textit{m} (plus \textit{n} \textit{p})) := \textbf{by} \textbf{rfl}
          _ = plus (.suc \textit{m}) (plus \textit{n} \textit{p}) := \textbf{by} \textbf{rfl}
\end{LeanVerbatim}
Let's unpack this code. The signature states that we are defining a theorem,
\LeanVerb|plus_assoc|, which states a proposition:\begin{verbatim}
\PLFAForall{} (m n p : \PLFANat{}) -> (m + n) + p = m + (n + p)
\end{verbatim}
The upside down A is pronounced "for all", and the proposition asserts that for
all natural numbers \LeanVerb|m|, \LeanVerb|n|, and \LeanVerb|p| the equation \LeanVerb|(m + n) + p = m + (n + p)|
holds. Evidence for the proposition is a function that accepts three natural
numbers, binds them to \LeanVerb|m|, \LeanVerb|n|, and \LeanVerb|p|, and returns evidence for the
corresponding instance of the equation.For the base case, we must show:\begin{verbatim}
plus (plus .zero n) p = plus .zero (plus n p)
\end{verbatim}
Simplifying both sides with the base case of addition yields the equation:\begin{verbatim}
plus n p = plus n p
\end{verbatim}
This holds trivially. Reading the chain of equations in the base case of the
proof, the top and bottom of the chain match the two sides of the equation to
be shown, and reading down from the top and up from the bottom takes us to
\LeanVerb|n + p| in the middle. No justification other than simplification is required.

\par

\subsection{More on rewriting}\label{PLFA--in--Lean----Part--1----Induction-005F-005F-005F--Proof--by--Induction----More--on--rewriting}

\par
\textbf{NOTE:} The last step may look a little backwards at first. Right before the end
of the inductive case (after rewriting using the inductive hypothesis \LeanVerb|ih|)
we are left with the term:

\par
\par\noindent \begin{verbatim}
.suc (plus m (plus n p))
\end{verbatim}
but the right-hand side of the theorem wants to look like:\begin{verbatim}
plus (.suc m) (plus n p)
\end{verbatim}
The important point is that the defining equation for \LeanVerb|plus| goes the other
way when it computes:\begin{verbatim}
plus (.suc m) k  -->  .suc (plus m k)
\end{verbatim}
Here is the definition of \LeanVerb|plus| recalled:\begin{verbatim}
def plus : \PLFANat{} -> \PLFANat{} -> \PLFANat{}
    | .zero , n    => n
    | (.suc m) , n => .suc (plus m n)
\end{verbatim}
So the term that directly matches the second clause of plus is not
\LeanVerb|.suc (plus m (plus n p))|. The term that directly matches is:\begin{verbatim}
plus (.suc m) (plus n p)
\end{verbatim}
and computing it gives:\begin{verbatim}
.suc (plus m (plus n p))
\end{verbatim}
So in the written proof, when we move from:\begin{verbatim}
.suc (plus m (plus n p))
\end{verbatim}
to:\begin{verbatim}
plus (.suc m) (plus n p)
\end{verbatim}
we are visually moving from the computed/unfolded form back to the folded-up
form. From a human reader point of view, it is fine to think of this as
"folding the definition of plus back up". But \LeanVerb|rfl| is not really doing a
directional rewrite here. It is checking that both sides compute to the same
expression:\begin{verbatim}
plus (.suc m) (plus n p)
-- computes to
.suc (plus m (plus n p))
\end{verbatim}
In other words, both sides have the same normal form, so Lean accepts the step
by \LeanVerb|rfl|. This is the same kind of silent definitional-equality step that Agda
writes explicitly as \LeanVerb|\PLFAEquiv{}\PLFASymbol{⟨}\PLFASymbol{⟩}|.

\par

\subsection{Second proof: commutativity}\label{PLFA--in--Lean----Part--1----Induction-005F-005F-005F--Proof--by--Induction----Second--proof-005F-005F-005F--commutativity}

\par
Another import property of addition is that it is \emph{commutative}, that is, that
the the order of the operands does not matter:
\LeanVerb|plus m n = plus n m|

\par
\par\noindent The proof requires that we first demonstrate two lemmas.

\par

\subsubsection{The first lemma}\label{PLFA--in--Lean----Part--1----Induction-005F-005F-005F--Proof--by--Induction----Second--proof-005F-005F-005F--commutativity----The--first--lemma}

\par
The base case of the definition of addition states that zero is a left identity.

\par
\par\noindent \begin{verbatim}
plus .zero n = n
\end{verbatim}
Our first lemma states that zero is also a right identity:\begin{verbatim}
plus m .zero = m
\end{verbatim}
Here is the lemma's statement and proof:\begin{LeanVerbatim}
\textbf{theorem} plusIdentityRight : \PLFAForall{} (\textit{m} : \PLFANat{}),  \textit{m} + .zero = \textit{m} := \textbf{by}
    \textbf{intro} \textit{m}
    \textbf{induction} \textit{m} \textbf{with}
    \PLFAComment{-- goal: plus \PLFANat{}.zero \PLFANat{}.zero = \PLFANat{}.zero}
    \symbol{124} zero =>
        \textbf{calc}
            plus \PLFANat{}.zero \PLFANat{}.zero
            = \PLFANat{}.zero            := \textbf{by} \textbf{rfl}
    \PLFAComment{-- goal: plus (\PLFANat{}.suc m) \PLFANat{}.zero = \PLFANat{}.suc m}
    \PLFAComment{--       ih : plus m \PLFANat{}.zero = m}
    \symbol{124} suc \textit{m} \textit{ih} =>
        \textbf{calc}
            plus (\PLFANat{}.suc \textit{m}) \PLFANat{}.zero
            = .suc (plus \textit{m} .zero) := \textbf{by} \textbf{rfl}
          _ = .suc \textit{m}              := \textbf{by} \textbf{exact} congrArg \PLFANat{}.suc \textit{ih}
\end{LeanVerbatim}
The signature states that we are defining the identifier \LeanVerb|plusIdentityRight|
which provides evidence for the proposition:\begin{verbatim}
\PLFAForall{} (m : \PLFANat{}),  m + .zero = m
\end{verbatim}
evidence for the proposition is a function that accepts a natural number, binds
it to \LeanVerb|m|, and returns evidence for the corresponding instance of the equation.
The proof uses lean4's induction tactic on \LeanVerb|m|.

\par
For the base case, we must show:

\par
\par\noindent \begin{verbatim}
plus .zero .zero = .zero
\end{verbatim}
simplifying with the base case of addition, this is straightforward.

\par
For the inductive case, we must show:

\par
\par\noindent \begin{verbatim}
plus (.suc m) .zero = .suc m
\end{verbatim}
Simplifying both sides with the inductive case of addition yields the equation:\begin{verbatim}
.suc (.plus m zero) = .suc m
\end{verbatim}
This in turn follows by prefacing \LeanVerb|.suc| to both sides of the induction hypothesis:\begin{verbatim}
plus m .zero = m
\end{verbatim}
Reading the chain of equations down from the top and up from the bottom takes us to
the simplified equation above. The remaining equation has the justification:\begin{verbatim}
exact congrArg \PLFANat{}.suc ih
\end{verbatim}
where \LeanVerb|exact| is a tactic that automatically closes the goal in the event that the
shape of the equation that results from \LeanVerb|congrArg N.suc ih| is automatically dischargeable...
\LeanVerb|congArg N.suc ih| has the effect of tacking the successor (\LeanVerb|.suc|) on both the
left hand side and right hand side of the inductive hypothesis arg... here is the sig
for congruence arg function \LeanVerb|congArg|:\begin{verbatim}
congArg: \PLFAForall{} {α β : Sort} ->
         \PLFAForall{} {a1 a2 : α} ->
         \PLFAForall{} (f : α -> β) ->
         \PLFAForall{} (h : a1 = a2) -> f a1 = f a2 := ...
\end{verbatim}

\par

\subsubsection{The second lemma}\label{PLFA--in--Lean----Part--1----Induction-005F-005F-005F--Proof--by--Induction----Second--proof-005F-005F-005F--commutativity----The--second--lemma}

\par
The inductive case of the definition of addition pushes \LeanVerb|.suc| on the first
argument to the outside:

\par
\par\noindent \begin{verbatim}
plus (.suc m) n = .suc (plus m n)
\end{verbatim}
Our second lemma does the same for \LeanVerb|.suc| on the second argument:\begin{verbatim}
plus m (.suc n) = .suc (plus m n)
\end{verbatim}
Here is the lemma's statement and proof:\begin{LeanVerbatim}
\textbf{theorem} plusSuc : \PLFAForall{} (\textit{m} \textit{n} : \PLFANat{}), plus \textit{m} (.suc \textit{n}) = .suc (plus \textit{m} \textit{n}) := \textbf{by}
    \textbf{intro} \textit{m} \textit{n}
    \textbf{induction} \textit{m} \textbf{with}
    \PLFAComment{-- goal: plus \PLFANat{}.zero (\PLFANat{}.suc n) = \PLFANat{}.suc (plus \PLFANat{}.zero n)}
    \symbol{124} zero =>
        \textbf{calc}
            plus \PLFANat{}.zero (\PLFANat{}.suc \textit{n})
            = .suc \textit{n}                  := \textbf{by} \textbf{rfl}
          _ = .suc (plus .zero \textit{n})     := \textbf{by} \textbf{rfl} \PLFAComment{-- rewrite using first defining eq of plus}
    \PLFAComment{-- goal: plus (\PLFANat{}.suc m) (\PLFANat{}.suc n) = \PLFANat{}.suc (plus (\PLFANat{}.suc m) n)}
    \PLFAComment{-- ih  : plus m (\PLFANat{}.suc n) = \PLFANat{}.suc (plus m n)}
    \symbol{124} suc \textit{m} \textit{ih} =>
        \textbf{calc}
            plus (.suc \textit{m}) (.suc \textit{n})
            = .suc (plus \textit{m} (.suc \textit{n})) := \textbf{by}
                \textbf{rfl}
          _ = \PLFANat{}.suc (\PLFANat{}.suc (plus \textit{m} \textit{n})) := \textbf{by}
                \PLFAComment{-- useful trick w/ hover on h for}
                \PLFAComment{-- examining type shape ("what if"?)}
                \textbf{have} \textit{h} := congrArg \PLFANat{}.suc \textit{ih}
                \textbf{exact} \textit{h}
          _ = \PLFANat{}.suc (plus (.suc \textit{m}) \textit{n}) := \textbf{by}
                \textbf{rfl}
\end{LeanVerbatim}
The signature states that we are defining the theorem \LeanVerb|plusSuc|, which provides
evidence for the proposition:\begin{verbatim}
\PLFAForall{} (m n : \PLFANat{}), plus m (.suc n) = .suc (plus m n)
\end{verbatim}
Evidence for this proposition is a function that accepts two natural numbers,
binds them to m and n, and returns evidence for the corresponding instance of
the equation.

\par
\par\noindent The proof is by induction on m. For the base case, we must show:\begin{verbatim}
plus .zero (.suc n) = .suc (plus .zero n)
\end{verbatim}
Simplifying both sides with the base case of addition, this is straightforward:\begin{verbatim}
| zero =>
    calc
        plus \PLFANat{}.zero (\PLFANat{}.suc n)
        = .suc n              := by rfl
      _ = .suc (plus .zero n) := by rfl
\end{verbatim}

\par
\par\noindent For the inductive case, we must show:\begin{verbatim}
plus (.suc m) (.suc n) = .suc (plus (.suc m) n)
\end{verbatim}
Simplifying the left-hand side with the inductive case of addition gives:\begin{verbatim}
.suc (plus m (.suc n))
\end{verbatim}
Simplifying the desired right-hand side also exposes:\begin{verbatim}
.suc (.suc (plus m n))
\end{verbatim}
So the proof passes through the middle expression:\begin{verbatim}
.suc (.suc (plus m n))
\end{verbatim}
The induction hypothesis is:\begin{verbatim}
ih : plus m (.suc n) = .suc (plus m n)
\end{verbatim}
Prefacing \LeanVerb|.suc| to both sides of the induction hypothesis gives:\begin{verbatim}
congrArg \PLFANat{}.suc ih :
  \PLFANat{}.suc (plus m (.suc n)) = \PLFANat{}.suc (\PLFANat{}.suc (plus m n))
\end{verbatim}
This is the central step of the calculation. Reading the chain from the
top down and from the bottom up takes us to the simplified equation in the
middle:\begin{verbatim}
| suc m ih =>
    calc
        plus (.suc m) (.suc n)
        = .suc (plus m (.suc n)) := by rfl
      _ = \PLFANat{}.suc (\PLFANat{}.suc (plus m n)) := by
            exact congrArg \PLFANat{}.suc ih
      _ = \PLFANat{}.suc (plus (.suc m) n) := by rfl
\end{verbatim}
Here, \LeanVerb|ih| is the induction hypothesis, and \LeanVerb|congrArg \PLFANat{}.suc ih| is the Lean
analogue of applying congruence with suc: it places \LeanVerb|\PLFANat{}.suc| around both sides
of the equality supplied by \LeanVerb|ih|.

\par
The first and last steps are justified by \LeanVerb|rfl|, because both are definitional
equalities following from the recursive definition of \LeanVerb|plus|.

\par

\subsubsection{The proposition}\label{PLFA--in--Lean----Part--1----Induction-005F-005F-005F--Proof--by--Induction----Second--proof-005F-005F-005F--commutativity----The--proposition}

\par
\par\noindent Finally, here is our proposition's statement and proof:\begin{LeanVerbatim}
\textbf{theorem} plusComm : \PLFAForall{} (\textit{m} \textit{n} : \PLFANat{}), \textit{m} + \textit{n} = \textit{n} + \textit{m} := \textbf{by}
    \textbf{intro} \textit{m} \textit{n}
    \textbf{induction} \textit{n} \textbf{with}
    \PLFAComment{-- goal: plus m \PLFANat{}.zero = plus \PLFANat{}.zero m}
    \symbol{124} zero =>
        \textbf{calc}
            plus \textit{m} .zero
          = \textit{m}            := \textbf{by} \textbf{exact} (plusIdentityRight \textit{m})
        _ = plus .zero \textit{m} := \textbf{by} \textbf{rfl}
    \symbol{124} suc \textit{n} \textit{ih} =>
        \PLFAComment{-- plus m (\PLFANat{}.suc n) = plus (\PLFANat{}.suc n) m}
        \PLFAComment{-- ih : plus m n = plus n m}
        \textbf{calc}
            plus \textit{m} (\PLFANat{}.suc \textit{n})
          = .suc (plus \textit{m} \textit{n}) := \textbf{by} \textbf{exact} plusSuc \textit{m} \textit{n}
        _ = .suc (plus \textit{n} \textit{m}) := \textbf{by}
            \PLFAComment{-- h : \PLFANat{}.suc (plus m n) = \PLFANat{}.suc (plus n m)}
            \textbf{have} \textit{h} := congrArg \PLFANat{}.suc \textit{ih}
            \textbf{exact} \textit{h}
        _ = plus (.suc \textit{n}) \textit{m} := \textbf{by} \textbf{rfl}
\end{LeanVerbatim}
The first line states that we are defining the identifier \LeanVerb|plusComm| which
provides evidence for the proposition:\begin{verbatim}
\PLFAForall{} (m n : \PLFANat{}), plus m n = plus n m
\end{verbatim}
Evidence for the proposition is a function that accepts two natural numbers,
binds them to \LeanVerb|m| and \LeanVerb|n|, and returns evidence for the corresponding instance
of the equation. The proof is by \LeanVerb|induction| on \LeanVerb|n|. (Not on \LeanVerb|m| this time!)

\par
For the base case, we must show:

\par
\par\noindent \begin{verbatim}
plus m zero = plus zero m
\end{verbatim}
Simplifying both sides with the base case of addition yields the equation:\begin{verbatim}
plus m .zero = m
\end{verbatim}
The remaining equation has the justification \LeanVerb|(plusIdentityRight m)|, which
invokes the first lemma.For the inductive case, we must show:\begin{verbatim}
plus m (.suc n) = plus (.suc n) m
\end{verbatim}
Simplifying both sides with the inductive case of addition yields the equation:\begin{verbatim}
plus m (.suc n) = .suc (plus n m)
\end{verbatim}
We show this in two steps. First, we have:\begin{verbatim}
plus m (.suc n) = .suc (plus m n)
\end{verbatim}
which is justified by the second lemma, \LeanVerb|(plusSuc m n)|. Then we have\begin{verbatim}
.suc (plus m n) = .suc (plus n m)
\end{verbatim}
which is justified by congruence and the induction hypothesis,
\LeanVerb|congrArg \PLFANat{}.suc ih|. This completes the proof.
Lean4 requires that identifiers are defined before they are used, so we must
present the lemmas before the main proposition, as we have done above.
In practice, one will often attempt to prove the main proposition first, and
the equations required to do so will suggest what lemmas to prove.

\par

\subsubsection{Our first corollary: rearranging}\label{PLFA--in--Lean----Part--1----Induction-005F-005F-005F--Proof--by--Induction----Second--proof-005F-005F-005F--commutativity----Our--first--corollary-005F-005F-005F--rearranging}

\par
We can apply associativity to rearrange parentheses however we like. Here is an
example:

\par
\begin{verbatim}
theorem plusRearrange : \PLFAForall{} (m n p q : \PLFANat{}),
    plus (plus m n) (plus p q) = plus (plus m (plus n p)) q := by
    intro m n p q
    calc
        plus (plus m n) (plus p q)
        = plus (plus (plus m n) p) q := by
          symm
          have h := plusAssoc (plus m n) p q
          exact h
      _ = plus (plus m (plus n p)) q := by
          have h  := plusAssoc m n p
          -- tacks a plus σ q onto the rhs of whatever shape σ is
          have h' := congrArg (\PLFASymbol{λ} σ => plus σ q) h
          exact h'
\end{verbatim}

\par
\par\noindent No induction is required, we simply apply associativity twice.
A few points are worthy of note.First, addition associates to the left, so m + (n + p) + q stands for (m + (n + p)) + q.
Second, we use sym to interchange the sides of an equation. Proposition +-assoc (m + n) p q shifts parentheses from left to right:\begin{verbatim}
plus ((plus m n) p) q = plus (plus m n) (plus p q)
\end{verbatim}
To shift them the other way, we use sym (\LeanVerb|plusAssoc (plus m n) p q|):\begin{verbatim}
(m + n) + (p + q) \PLFAEquiv{} ((m + n) + p) + q
\end{verbatim}
In general, if \LeanVerb|e| provides evidence for \LeanVerb|x = y| then \LeanVerb|symm| tactic
(which must be invoked on its own line in a \LeanVerb|by| block) provides evidence for
\LeanVerb|y = x|.Third, Lean4 allows us to "tack" a term onto the rhs of both sides of an
equality by writing \LeanVerb|congrArg (\PLFASymbol{λ} σ => plus σ q) h|.. so this:\begin{verbatim}
plus (plus m n) p  =  plus m (plus n p)
\end{verbatim}
can be rewritten into the equation:\begin{verbatim}
plus (plus (plus m n) p) q  =  plus (plus m (plus n p)) q
\end{verbatim}

\par

\subsubsection{Creation, one last time}\label{PLFA--in--Lean----Part--1----Induction-005F-005F-005F--Proof--by--Induction----Second--proof-005F-005F-005F--commutativity----Creation-005F-005F-005F--one--last--time}

\par
Returning to the proof of associativity, it may be helpful to view the
inductive proof (or, equivalently, the recursive definition) as a creation
story. This time we are concerned with judgments asserting associativity:

\par
\par\noindent \begin{verbatim}
-- in the beginning, we know nothing about associativity
\end{verbatim}
Now, we apply the rules to all the judgments we know about. The base case tells us that (zero + n) + p \PLFAEquiv{} zero + (n + p) for every natural n and p. The inductive case tells us that if (m + n) + p \PLFAEquiv{} m + (n + p) (on the day before today) then (suc m + n) + p \PLFAEquiv{} suc m + (n + p) (today). We didn’t know any judgments about associativity before today, so that rule doesn’t give us any new judgments:\begin{verbatim}
-- on the first day, we know about associativity of 0
plus (plus 0 0) 0 = plus 0 (plus 0 0) ... plus (plus 0 4) 5 = plus 0 (plus 4 5)
\end{verbatim}
Then we repeat the process, so on the next day we know about all the judgments
from the day before, plus any judgments added by the rules. The base case tells
us nothing new, but now the inductive case adds more judgments:\begin{verbatim}
-- on the second day, we know about associativity of 0 and 1
plus (plus 0 0) 0 = plus 0 (plus 0 0) ... plus (plus 0 4) 5 = plus 0 (plus 4 5)
plus (plus 1 0) 0 = plus 1 (plus 0 0) ... plus (plus 1 4) 5 = plus 1 (plus 4 5)   ...
\end{verbatim}
And we repeat the process again:\begin{verbatim}
-- on the third day, we know about associativity of 0, 1, and 2
plus (plus 0 0) 0 = plus 0 (plus 0 0) ... plus (plus 0 4) 5 = plus 0 (plus 4 5)
plus (plus 1 0) 0 = plus 1 (plus 0 0) ... plus (plus 1 4) 5 = plus 1 (plus 4 5)   ...
plus (plus 2 0) 0 = plus 2 (plus 0 0) ... plus (plus 2 4) 5 = plus 2 (plus 4 5)   ...
\end{verbatim}
You've got the hang of it by now:\begin{verbatim}
-- On the fourth day, we know about associativity of 0, 1, 2, and 3.
plus (plus 0 0) 0 = plus 0 (plus 0 0) ... plus (plus 0 4) 5 = plus 0 (plus 4 5)
plus (plus 1 0) 0 = plus 1 (plus 0 0) ... plus (plus 1 4) 5 = plus 1 (plus 4 5)   ...
plus (plus 2 0) 0 = plus 2 (plus 0 0) ... plus (plus 2 4) 5 = plus 2 (plus 4 5)   ...
plus (plus 3 0) 0 = plus 3 (plus 0 0) ... plus (plus 3 4) 5 = plus 3 (plus 4 5)   ...
\end{verbatim}
The process continues. On the m'th day we will know all the judgments where the
first number is less than m.There is also a completely finite approach to generating the same equations,
which is left as an exercise for the reader.

\par

\subsubsection{Exercise \texttt{plusSwap} (recommended)}\label{PLFA--in--Lean----Part--1----Induction-005F-005F-005F--Proof--by--Induction----Second--proof-005F-005F-005F--commutativity----Exercise----plusSwap-----005FLPAR-005Frecommended-005FRPAR-005F}

\par
\par\noindent Show that\begin{verbatim}
plus m (plus n p) = plus n (plus m p)
\end{verbatim}
for all naturals \LeanVerb|m|, \LeanVerb|n|, and \LeanVerb|p|. No induction is needed, just apply the
previous results which show addition is associative and commutative.\textbf{potential sol.}todo

\par

\end{document}
